% Generated draft: requires mathematical and source review.
\documentclass[11pt]{article}
\usepackage[margin=25mm]{geometry}
\usepackage{fontspec}
\setmainfont{Latin Modern Roman}
\usepackage{amsmath,amssymb,mathtools,unicode-math}
\setmathfont{Latin Modern Math}
\usepackage{enumitem,needspace,xurl,hyperref}
\setlength{\parindent}{0pt}
\setlength{\parskip}{5pt}
\setlength{\emergencystretch}{3em}
\newcommand{\N}{\mathbb{N}}
\newcommand{\Z}{\mathbb{Z}}
\newcommand{\Q}{\mathbb{Q}}
\newcommand{\R}{\mathbb{R}}
\newcommand{\C}{\mathbb{C}}
\newcommand{\F}{\mathbb{F}}
\newcommand{\A}{\mathbb{A}}
\newcommand{\PP}{\mathbb{P}}
\newcommand{\E}{\mathbb{E}}
\newcommand{\eps}{\varepsilon}
\newcommand{\dd}{\,\mathrm{d}}
\begin{document}
\section*{UnsolvedMath ID: 10400167}
\textbf{Source problem code:} catalogue entry
\subsection{Definitions and mathematical statement}
A fusion-rule algebra has a distinguished finite basis of labels and nonnegative integer structure constants \(N_{ij}^k\) describing products. In a categorical realization, \(6j\)-symbols are matrix coefficients of associativity transformations on the corresponding multiplicity spaces, subject to coherence identities. They depend on basis choices; the extra structures required for a particular TQFT construction must also be supplied.

Suppose two admissible fusion/\(6j\) data sets \(D_1,D_2\) produce TQFTs \(Z_1,Z_2\), and these theories are isomorphic in the intended sense. Determine what relation between the two sets of \(6j\)-symbols follows, or equivalently characterize the fibers of the construction from admissible data to TQFT isomorphism classes.

The data need not have the same fusion algebra, label set or multiplicity dimensions. A basis-gauge or relabeling change is therefore not silently asserted to be the only possible relation. The source does not specify the state-sum/construction, coefficient field, pivotal/spherical/unitary assumptions, tangential structure, extendedness or exact TQFT isomorphism. A closed-manifold invariant, an unextended theory and a theory including boundary data may retain different information. These choices are essential before a precise necessary relation or equivalence criterion can be stated.

The summary supplies only an open/triage attribution and no exact theorem. The 2002 question remains source information pending current primary-literature review, rather than a claim that all variants have one unresolved classification.
\subsection{Short English statement}
If two fusion-rule and \(6j\)-symbol constructions give isomorphic TQFTs, how are the input data related? The construction and level of TQFT structure must be fixed; a common label set or mere gauge equivalence is not assumed.
\subsection{Sources}
\begin{itemize}
\item Local registry: \url{data/unsolvedmath/problems.json}, permanent UnsolvedMath ID \(10400167\), source code catalogue entry; supplied statement, background, and research summary, updated 2026-08-06.
\item Ohtsuki, \emph{Problems on invariants of knots and 3-manifolds} (2002), Problem 9.3, attributed to Y. Kawahigashi, PDF page 138: \url{https://msp.org/gtm/2002/04/gtm-2002-04-024s.pdf}. Supplied PDF-font-structure extraction accessed 2026-07-29; provenance distinguishes 231 numbered items.
\item No exact reconstruction/equivalence theorem or modern resolution reference is supplied.
\end{itemize}
\end{document}
